\section{The Ceiling of the Ascent}

To have God: the tradition's name for what that would be is exact to the point of audacity. Not nearness, in some spatial metaphor the metaphysics already retired---Part I found God nearer to each thing than its own core, and the heart stayed restless. Not information, however elevated. Aquinas puts the requirement in terms of the mind's own hunger and follows it without flinching: ``man is not perfectly happy, so long as something remains for him to desire and seek''; and something \emph{does} remain, so long as the first cause is known only from its effects---known, that is, in the manner Part I's whole ascent knew it, as the unseen singer inferred from the song. ``If therefore the human intellect, knowing the essence of some created effect, knows no more of God than `that He is,'\,'' its perfection has not yet reached the source; the natural desire to know \emph{what} the cause is stands unfulfilled. ``Final and perfect happiness can consist in nothing else than the vision of the Divine Essence.''\endnote{\emph{Summa theologiae} I-II, q.~3, a.~8, corpus (\latin{ultima et perfecta beatitudo non potest esse nisi in visione divinae essentiae}); checked 2026-07-25, English at New Advent, Latin at Corpus Thomisticum. The corpus concludes that the intellect ``needs to reach the very Essence of the First Cause. And thus it will have its perfection through union with God as with that object, in which alone man's happiness consists.'' The article stands in a work of Christian theology, and its author holds on faith that the vision is in fact man's destiny; what the argument itself supplies is the conditional: nothing less would close the distance between the will's calibration and its possession. This essay uses it at that level.} Seeing, not concluding. The difference is the difference between reading a fine description of a face and turning around.

The same reasoning appears at the very front of the \emph{Summa}, wearing its most disarming clothes. ``There resides in every man a natural desire to know the cause of any effect which he sees; and thence arises wonder in men. But if the intellect of the rational creature could not reach so far as to the first cause of things, the natural desire would remain void.''\endnote{\emph{Summa theologiae} I, q.~12, a.~1, corpus (\latin{inest enim homini naturale desiderium cognoscendi causam, cum intuetur effectum; et ex hoc admiratio in hominibus consurgit\ldots{} remanebit inane desiderium naturae}); checked 2026-07-25 as above. The corpus concludes: ``Hence it must be absolutely granted that the blessed see the essence of God.'' How much the natural-desire argument proves, and in what sense the desire is ``natural,'' is one of the most contested questions in the Thomistic tradition; the appendix states the boundary this essay observes, and nothing below treats the desire as a claim on its own fulfillment.} Wonder---the child's why, dignified by a Greek name---is the intellect's form of the heart's restlessness, and it aims past every intermediate answer at the first answer, not as a proposition but as a presence. A universe in which that aim were structurally futile would be a universe with a hunger built in and no food in it anywhere. Aquinas declines to believe in such a universe. But watch precisely what he asserts and where he asserts it: writing as a Christian theologian, he concludes that the blessed \emph{do} see the essence of God. The philosophical thread alone delivers something more austere: that nothing short of such seeing could be the end the will is calibrated to, and that a creature is the sort of thing whose desire runs that far. Whether the desire will be met is another matter entirely---and here the ascent, for the third time in this series, runs into its own ceiling, this time with a force the first two arrests only rehearsed.

For the seeing in question is not available for the climbing. Aquinas is as blunt about this as about the desire: ``the created intellect cannot see the essence of God, unless God by His grace unites Himself to the created intellect, as an object made intelligible to it.''\endnote{\emph{Summa theologiae} I, q.~12, a.~4, corpus; checked 2026-07-25 as above. The reason is structural: knowledge is according to the mode of the knower, and the mode of every created knower is proportioned to created natures, whereas ``to know self-subsistent being is natural to the divine intellect alone.'' Compare I-II, q.~5, a.~5, corpus: perfect happiness ``surpasses the nature not only of man, but also of every creature,'' so that ``neither man, nor any creature, can attain final Happiness by his natural powers.''} And in the treatise on happiness: perfect beatitude ``surpasses the nature not only of man, but also of every creature''; neither man nor angel can reach it by natural powers. Weigh what this does to the picture. The audit of the last section aimed every human wanting at one address. This section's first argument showed that having the address means seeing, not inferring. And now the same analysis certifies that no creature can see---not by effort, not by genius, not by ten thousand years of purification, not by being promoted through every rank of finite excellence, because the gap is not a distance within the creaturely order but the border of that order itself. The creature is calibrated to an end it cannot attain. Its reach is built longer than its arms.

This is the ontological vertigo again, in its final and most intimate register. Part I found the creature suspended over its own existing; Part II found it indebted beyond settlement; here it hangs over its own \emph{wanting}: the deepest desire in it is for something no exertion of its own can procure, and the honest ascent ends not at a summit but at a locked ceiling, with the compass still pointing up. Recoil has its move ready, as always: call the calibration a cruelty---a hunger installed, on this reading, as a taunt. And as always, recoil has misread the same asymmetry one section too early. For Aquinas notices what kind of creature you would have to be for the ceiling to be an insult, and answers with a comparison of striking gentleness: a nature ``that can attain perfect good, although it needs help from without in order to attain it, is of more noble condition than a nature which cannot attain perfect good, but attains some imperfect good, although it need no help from without in order to attain it''---as a man curable by medicine is better off in the order of health than one whose constitution can manage only a lesser health, however self-sufficiently.\endnote{\emph{Summa theologiae} I-II, q.~5, a.~5, ad~2, citing Aristotle, \emph{De caelo} ii, 12; checked 2026-07-25 as above. The reply's comparison (``he is better disposed to health who can attain perfect health, albeit by means of medicine, than he who can attain but imperfect health, without the help of medicine'') is Aquinas's own; its application here---the open mouth as the shape of a creature made to be fed---is this essay's image.} The beasts have no locked ceiling; nothing in a sparrow reaches past the creaturely order, and nothing in a sparrow is disappointed. To carry a desire only God can satisfy is not the design flaw of the human creature. It is its rank. The open mouth is not the mark of the starving; it is the shape of a creature made to be fed---made, as Augustine said with his preposition, \emph{for}.

But a shape is not a meal, and here the philosophical thread of this series reaches the end it has been approaching since Part I named its own boundary at the Trinity. Whether the source will give what no creature can take is not a question any argument can settle, for the same reason no argument in Part I could have predicted creation: a gift, by definition, cannot be deduced. The ascent can establish the calibration, the insufficiency of every creature, the identity of the only sufficient end, and the impossibility of self-help. It can establish, in other words, exactly this: that the creature's situation is a question addressed to God---asked not by its words but by its structure, the way the open mouth of a nestling is a question. What philosophy cannot do is answer for Him. If an answer exists, it will have to come from the other side, and it will have to be, once more and finally, a gift. The remainder of this essay listens to what the tradition testifies that answer to be; and the seam between what has been argued and what is now received---marked in every part of this series---is hereby marked for the last time, at the exact place where the creature runs out of ladder.
